\documentclass[11pt,a4paper]{article}

% ── Packages ──────────────────────────────────────────────────────────────────
\usepackage{amsmath,amssymb,amsthm}
\usepackage{graphicx}
\usepackage{xcolor}
\usepackage{geometry}
\usepackage{booktabs}
\usepackage{tabularx}
\usepackage{adjustbox}
\usepackage{enumitem}
\usepackage{tcolorbox}
\usepackage{colortbl}
\usepackage{multirow}
\usepackage{float}
\usepackage{caption}

% ── hyperref (last) ──────────────────────────────────────────────────────────
\usepackage[
    colorlinks=true,
    linkcolor=blue!70!black,
    citecolor=darkgray,
    urlcolor=blue!70!black,
    bookmarks=true,
    bookmarksnumbered=true,
    unicode=true
]{hyperref}

% ── Geometry ─────────────────────────────────────────────────────────────────
\geometry{a4paper, top=2.5cm, bottom=2.5cm, left=2.5cm, right=2.5cm}

% ── Theorem environments ─────────────────────────────────────────────────────
\theoremstyle{definition}
\newtheorem{definition}{Definition}[section]
\newtheorem{theorem}{Theorem}[section]
\newtheorem{lemma}[theorem]{Lemma}
\newtheorem{proposition}{Proposition}[section]

% ── Custom colors ────────────────────────────────────────────────────────────
\definecolor{accentblue}{HTML}{1B3A5C}
\definecolor{accentorange}{HTML}{D4763C}
\definecolor{lightgray}{HTML}{F5F5F5}
\definecolor{medgray}{HTML}{E0E0E0}

% ── Listings ─────────────────────────────────────────────────────────────────
\usepackage{listings}
\lstset{
  basicstyle=\ttfamily\footnotesize,
  breaklines=true,
  frame=none,
  backgroundcolor=\color{lightgray},
  xleftmargin=1em,
  xrightmargin=1em
}

% ── Orphan/widow control ─────────────────────────────────────────────────────
\widowpenalty=10000
\clubpenalty=10000

% ── Figure path ──────────────────────────────────────────────────────────────
\graphicspath{{./paper_assets/}}

% ── PDF metadata ─────────────────────────────────────────────────────────────
\hypersetup{
    pdftitle={From Coding Theory to Harmony: How Mersenne and Fermat Primes Reveal the Number-Theoretic Structure of Musical Consonance},
    pdfauthor={Computational Musicology Research},
    pdfsubject={UBP System, Golay Code, Leech Lattice, Prime Numbers, Music Theory},
    pdfkeywords={computational musicology, Golay code, Leech lattice, Mersenne primes, Fermat primes, Jaccard similarity, consonance, UBP}
}

\begin{document}

% ═══════════════════════════════════════════════════════════════════════════
% NO TITLE PAGE IN .tex — Cover generated separately via HTML/Playwright
% Body starts with Table of Contents
% ═══════════════════════════════════════════════════════════════════════════

\tableofcontents
\newpage

% ═══════════════════════════════════════════════════════════════════════════
\section{Introduction}\label{sec:intro}
% ═══════════════════════════════════════════════════════════════════════════

The relationship between mathematical structure and musical harmony has fascinated scholars since Pythagoras first observed that consonant musical intervals correspond to simple integer ratios of string lengths. In the modern era, this question has been reframed in terms of coding theory, lattice geometry, and number theory: can the same algebraic structures that enable error-correcting codes also capture the perceptual phenomenon of consonance? This paper presents a ten-phase computational investigation conducted within the Universal Binary Principle (UBP) system---a software framework that implements the extended binary Golay code $[24,12,8]$, the Leech lattice $\Lambda_{24}$, and the Barnes-Wall lattice family $\Lambda_n$ for $n \in \{256, 512, 1024\}$.

The central question motivating this study is whether the UBP's mathematical infrastructure, originally designed for error correction and data integrity, can naturally encode musical relationships without ad hoc construction. Rather than building a music-theoretic model and fitting mathematical structures to it, we adopt the opposite approach: we begin with the mathematics and ask, at each stage, whether it captures anything recognizable as harmony. This methodology of ``honest mapping''---refusing to force or fabricate correlations---distinguishes this work from prior computational musicology efforts that select encodings to maximize fit.

The investigation proceeds through four arcs. The first arc (Phases I--V) establishes that the Circle-of-Fifths Gray encoding through the Golay code produces a strong positive correlation ($r = 0.8674$, $R^2 = 75.2\%$) between Hamming distance and interval consonance, but that this correlation is structurally shallow: it arises from a three-step chain (Circle-of-Fifths ordering $\to$ Gray code adjacency $\to$ Golay distance) that maps six consonance ranks onto only two of the Golay code's three possible inter-codeword distances. Chord-level differentiation proves impossible within this framework because the Golay code's three-distance constraint ($\{8, 12, 16\}$) provides insufficient resolution to separate even basic triad types.

The second arc (Phases VI--VII) demonstrates that dimensional expansion through the Leech lattice (24D) and Barnes-Wall lattices (256D, 512D, 1024D) preserves the interval correlation identically but never enriches it. The correlation $r = 0.8674$ persists unchanged across four orders of magnitude in dimension, while chord correlation remains flat at $r \approx -0.09$. This ``signal preservation without enrichment'' phenomenon is traced to the fact that all expansion is linear: higher-dimensional embeddings are linear projections that inherit the Golay code's distance distribution without adding new structure.

The third arc (Phase VIII) marks a turning point. When we abandon the coding-theoretic approach entirely and encode pitches using the prime exponent vectors of their just-intonation frequency ratios---that is, representing the perfect fifth $3/2$ as $(0, 1, 0)$ for the exponents of $(2, 3, 5)$---the interval correlation jumps to $r = +0.9613$, the strongest result in the entire study. This demonstrates unequivocally that the consonance signal lives in number theory (specifically, in the prime factorization of frequency ratios), not in coding theory. The same phase also discovers a remarkable duality between Mersenne primes ($2^p - 1$) and Fermat primes ($2^{2^k} + 1$) in modular arithmetic: both families, when reduced modulo $144 = 12^2$, produce residues that encode the perfect fifth and fourth---the two intervals that partition the octave.

The fourth arc (Phases IX--X) converts this prime-theoretic insight into a practical harmonic module. Phase IX shows that a four-dimensional vector of distances from each pitch class to the four prime residues $\{17, 31, 113, 127\}$ in mod-144 space achieves the best chord correlation in the entire study ($r = -0.8790$ via Euclidean distance, inverted because consonant chords spread apart while dissonant ones cluster). Phase X then introduces Jaccard similarity as the metric, and a proximity threshold of $\tau = 10$ yields the \emph{first positive chord differentiation} above $r = 0.5$ in the entire investigation: $r = +0.8244$ for the pairwise average Jaccard similarity. The mechanism is traced to the fact that $17 \equiv 5 \pmod{12}$ (the fourth, a Fermat prime) and $31 \equiv 7 \pmod{12}$ (the fifth, a Mersenne form), so that the four prime residues naturally partition mod-144 space at the tritone boundary, and consonant triads systematically span this boundary while dissonant clusters do not.

The principal contributions of this paper are: (1) a rigorous negative result showing that no coding-theoretic extension of the Golay code can differentiate chords, proven through exhaustive permutation search and dimensional scaling; (2) the discovery that $144 = 12^2$ is the ``correct'' modulus for connecting Mersenne/Fermat prime structure to the chromatic octave; (3) the identification of a Jaccard-based Prime-Layer Harmonic Module that achieves $r = +0.8244$ for chord differentiation---a positive, transposition-invariant, structurally explained mapping from prime number theory to musical harmony.

% ═══════════════════════════════════════════════════════════════════════════
\section{The UBP System Architecture}\label{sec:ubp}
% ═══════════════════════════════════════════════════════════════════════════

\subsection{The Extended Binary Golay Code}\label{sec:golay}

The UBP system's core is the extended binary Golay code $\mathcal{G}_{24}$, a $[24, 12, 8]$ linear error-correcting code over $\mathrm{GF}(2)$. Each of the $2^{12} = 4096$ codewords is a 24-bit binary vector. The code's minimum Hamming distance is $d_{\min} = 8$, meaning any two distinct codewords differ in at least 8 bit positions. The code is self-dual ($\mathcal{G}_{24} = \mathcal{G}_{24}^{\perp}$), doubly-even (all codewords have Hamming weight divisible by 4), and perfect in the sense that it achieves the sphere-packing bound in 24 dimensions.

A critical property for this study is the \emph{distance distribution} of $\mathcal{G}_{24}$. Despite having 4096 codewords, the set of possible inter-codeword Hamming distances contains only three values: $d \in \{8, 12, 16\}$. The weight enumerator polynomial is:
\begin{equation}\label{eq:weight_enum}
W(x) = 1 + 759x^8 + 2576x^{12} + 759x^{16} + x^{24}.
\end{equation}
This means that for any two codewords $\mathbf{c}_i, \mathbf{c}_j \in \mathcal{G}_{24}$, the Hamming distance $d_H(\mathbf{c}_i, \mathbf{c}_j) = w_H(\mathbf{c}_i \oplus \mathbf{c}_j)$ (since $\mathcal{G}_{24}$ is linear) can only be 8, 12, or 16, with the majority ($\approx 62.9\%$) of pairs at distance 12. Codewords of weight 8 are called \emph{octads} and possess exceptional symmetry properties. As we shall see, this three-value constraint is the fundamental obstacle to chord differentiation.

The UBP's encoding pipeline maps a 12-bit input seed to a 24-bit Golay codeword through a systematic generator matrix. For pitch encoding, each of the 12 pitch classes is assigned a 12-bit seed, which is then encoded to a 24-bit codeword. The choice of seed assignment (the mapping from pitch class to seed) is the primary encoding variable explored in this study.

\subsection{The Leech Lattice and Barnes-Wall Lattices}\label{sec:leech_bw}

The Leech lattice $\Lambda_{24}$ is a 24-dimensional Euclidean lattice that can be constructed from the Golay code. Each octad (weight-8 codeword) expands to 128 physical lattice points in $\mathbb{R}^{24}$ with coordinates in $\{0, \pm 2\}$. The lattice achieves the highest possible packing density in 24 dimensions, with a kissing number of 196560 (the number of nearest neighbors of any lattice point).

The Barnes-Wall lattice $\Lambda_{256}$ is a 256-dimensional lattice that can be recursively constructed from the Reed-Muller code $RM(1,8)$ and inherits structure from the Golay code. The UBP system implements Barnes-Wall lattices at dimensions 256, 512, and 1024. These lattices provide an ordered hierarchy of increasing dimensional embedding while maintaining the algebraic connection to the Golay code.

Both lattices are used in this study as ``dimensional expansion'' tools: we embed pitch class representations into higher-dimensional spaces to test whether additional dimensions provide finer distance resolution for harmonic differentiation. As Section~\ref{sec:expansion} demonstrates, they do not.

\subsection{Consonance Model}\label{sec:consonance}

Throughout this study, we use a standard six-level consonance ranking for the 11 chromatic intervals (Table~\ref{tab:consonance}). This ranking follows the traditional music-theoretic classification into perfect consonances (unison, octave, fifth, fourth), imperfect consonances (major and minor thirds, major and minor sixths), and dissonances (seconds, sevenths, tritone). Lower rank values indicate greater consonance.

\begin{table}[htbp]
\centering
\caption{Consonance ranking used throughout this study.}\label{tab:consonance}
\begin{tabular*}{\columnwidth}{@{\extracolsep{\fill}} c l l c}
\toprule
Rank & Interval & Abbreviation & Semitones \\
\midrule
1 & Unison / Octave & P1/P8 & 0 / 12 \\
2 & Perfect Fifth & P5 & 7 \\
3 & Perfect Fourth & P4 & 5 \\
3 & Major Third & M3 & 4 \\
3 & Major Sixth & M6 & 9 \\
4 & Minor Third & m3 & 3 \\
4 & Minor Sixth & m6 & 8 \\
4 & Minor Seventh & m7 & 10 \\
5 & Major Second & M2 & 2 \\
5 & Major Seventh & M7 & 11 \\
6 & Minor Second & m2 & 1 \\
6 & Tritone & TT & 6 \\
\bottomrule
\end{tabular*}
\end{table}

% ═══════════════════════════════════════════════════════════════════════════
\section{Pitch Encoding and Interval Dynamics}\label{sec:phases12}
% ═══════════════════════════════════════════════════════════════════════════

\subsection{Phase I: Three Encoding Strategies}\label{sec:phase1}

In Phase~I, we encoded the 12 pitch classes $\{0, 1, \ldots, 11\}$ into 24-bit Golay codewords using three distinct 12-bit seed assignments, then measured the Hamming distance between all $\binom{12}{2} = 66$ unordered pitch pairs and correlated these distances with the consonance ranks from Table~\ref{tab:consonance}. The three seed assignments were:

\begin{enumerate}[leftmargin=2em]
\item \textbf{Chromatic Gray}: Pitch class $i$ is encoded as the 12-bit binary-reflected Gray code of $i$. This assigns consecutive pitch classes to consecutive Gray code positions.
\item \textbf{Circle-of-Fifths (CoF) Gray}: Pitch classes are ordered by their position on the circle of fifths ($C, G, D, A, E, B, F\#, C\#/D\flat, A\flat, E\flat, B\flat, F$), and the position index (0--11) is Gray-encoded. This ordering reflects the most fundamental relationship in tonal harmony.
\item \textbf{One-Hot}: Each pitch class $i$ is encoded as a 12-bit vector with a single 1 at position $i$. This serves as a maximally sparse baseline with no structural assumptions.
\end{enumerate}

\subsection{Phase II: Interval Correlation Results}\label{sec:phase2}

The results, summarized in Figure~\ref{fig:encoding_comparison}, reveal a dramatic dependence on the choice of seed assignment. The Circle-of-Fifths Gray encoding achieves $r = +0.8674$ ($R^2 = 75.2\%$), a strong positive correlation indicating that more consonant intervals are mapped to smaller Hamming distances in Golay space. By contrast, the Chromatic Gray encoding achieves only $r = +0.2947$, and the One-Hot encoding produces $r = 0.0000$ (zero correlation, since all one-hot codewords lie at the same Hamming distance from each other). A direct interval encoding---where the interval width itself (0--11 semitones) is Gray-coded and Golay-encoded---yields $r = +0.52$, intermediate between the two orderings.

\begin{figure}[htbp]
\centering
\includegraphics[width=\columnwidth, height=0.35\textheight, keepaspectratio]{fig1_encoding_comparison.png}
\caption{Pearson correlation $r$ between consonance rank and Hamming distance for five encoding strategies. The Circle-of-Fifths Gray encoding through the Golay code achieves the strongest correlation at $r = 0.8674$.}\label{fig:encoding_comparison}
\end{figure}

The superiority of the Circle-of-Fifths ordering can be understood structurally. Gray code ensures that consecutive positions differ in exactly one bit, so the Circle-of-Fifths ordering means that musically adjacent pitches (separated by a perfect fifth) produce seeds differing in one bit. The Golay code's generator matrix then amplifies this 1-bit seed difference into a consistent 8-bit (octad) codeword difference for the most consonant intervals. The key insight is that the Golay code acts as a \emph{distance amplifier}: it maps the 1-bit Gray adjacency between fifth-related pitches onto the minimum Hamming distance of 8, while larger musical intervals (which require multiple CoF steps, and hence multiple bit flips) map onto larger Hamming distances.

However, as we show in Section~\ref{sec:chord_problem}, this correlation, while genuine, is structurally too coarse to extend from intervals to chords.

\subsection{Phase II--C: Gray Code Is Essential}

A critical control experiment demonstrated that the Gray code transformation is essential. When we use the Circle-of-Fifths ordering with \emph{raw binary} (not Gray code) seeds, the correlation drops from $r = 0.8674$ to $r = 0.18$. This confirms that the correlation depends on the specific distance structure imposed by the Gray code, not merely on the Circle-of-Fifths ordering itself. Gray code ensures that each CoF step changes exactly one bit, creating a smooth distance gradient that the Golay code then preserves and amplifies.

% ═══════════════════════════════════════════════════════════════════════════
\section{The Chord Problem and the Three-Distance Ceiling}\label{sec:chord_problem}
% ═══════════════════════════════════════════════════════════════════════════

\subsection{Phase III: Chordal XOR Synthesis}

In Phase~III, we extended the investigation from intervals to chords by representing a chord as the XOR (addition in $\mathrm{GF}(2)^{24}$) of its constituent pitch codewords. For a chord $\{p_1, p_2, \ldots, p_k\}$, the chord codeword is $\mathbf{c}_{\text{chord}} = \mathbf{c}_{p_1} \oplus \mathbf{c}_{p_2} \oplus \cdots \oplus \mathbf{c}_{p_k}$. Since the Golay code is linear, this XOR is also a valid codeword.

The result was complete failure: \emph{all} triads---major, minor, diminished, augmented, and dissonant clusters---produced XOR codewords with Hamming weight 8 (octads). Every 3-note combination of the 12 CoF-Gray encoded pitches XORs to an octad, making chord-level differentiation impossible within the Golay code's linear structure. The 39 XOR-closed 4-note chords are precisely the 3-dimensional flats in $\mathrm{GF}(2)^{12}$, a purely linear-algebraic artifact with no particular musical significance.

\subsection{Phase IV: Why the Circle-of-Fifths Works}

To understand the $r = 0.8674$ result, we conducted an exhaustive search over 10,000 random permutations of pitch-to-seed assignments (with Gray coding). The Circle-of-Fifths ordering ranked approximately \#1637 out of 50,000 (top 3.3\%), confirming it is genuinely good but not uniquely so. The Tritone Split ordering achieved $r = -0.89$ (strong but inverted), demonstrating that the Golay code can capture interval structure in either direction depending on the seed assignment.

A Monte Carlo analysis revealed that random 12-bit vectors provide an average of 8.7 distinct inter-codeword distances---more than enough for harmonic differentiation. The problem is not the bit space but the Golay code itself, which collapses these to only three distances.

\subsection{The Fundamental Obstacle}\label{sec:ceiling}

Figure~\ref{fig:golay_ceiling} illustrates the fundamental problem. The Golay code's weight enumerator (Equation~\ref{eq:weight_enum}) provides only three non-zero inter-codeword distances: $\{8, 12, 16\}$. In practice, only two of these appear for the 12 pitch codewords (distances 8 and 12). Six consonance ranks must be mapped onto two distance values, which is information-theoretically insufficient for fine-grained harmonic differentiation.

\begin{figure}[htbp]
\centering
\includegraphics[width=\columnwidth, height=0.35\textheight, keepaspectratio]{fig2_golay_ceiling.png}
\caption{The Golay $[24,12,8]$ three-distance ceiling. Left: the three possible inter-codeword distances (8, 12, 16) with their theoretical proportions. Right: the six consonance ranks collapse onto only two effective distance values (8 and 12), illustrating the information loss that prevents chord differentiation.}\label{fig:golay_ceiling}
\end{figure}

Formally, to separate six consonance ranks at the interval level, a code requires at least six distinct distance values. The Golay code provides three, of which only two appear for pitch codewords. For chord differentiation---separating at least five triad types (major, minor, diminished, augmented, dissonant)---the combinatorial requirement is even more severe, as each chord type is characterized by a multiset of three pairwise interval distances. The Golay code's constraint makes this fundamentally impossible: \emph{no metric defined on Golay codewords can differentiate chord types beyond what the three-distance constraint permits}.

% ═══════════════════════════════════════════════════════════════════════════
\section{Dimensional Expansion: Leech and Barnes-Wall}\label{sec:expansion}
% ═══════════════════════════════════════════════════════════════════════════

\subsection{Phase VI: Leech Lattice and Barnes-Wall 256D}

Given the Golay code's three-distance ceiling, a natural question arises: can higher-dimensional lattice embeddings provide finer distance resolution? The Leech lattice $\Lambda_{24}$ expands each octad (weight-8 Golay codeword) into 128 physical points in $\mathbb{R}^{24}$, while the Barnes-Wall lattice $\Lambda_{256}$ provides a 256-dimensional embedding.

The results, shown in Figure~\ref{fig:dimensional_hierarchy}, are striking. The Leech lattice's Euclidean distances between pitch representations destroy the interval correlation entirely ($r = 0.13$). This occurs because the 128-point symmetric expansion of each octad washes out the binary distance signal: when each pitch becomes a cloud of 128 points, the pairwise Euclidean distances between clouds converge to a narrow distribution, losing the two-bucket structure that drove the $r = 0.8674$ result.

\begin{figure}[htbp]
\centering
\includegraphics[width=\columnwidth, height=0.35\textheight, keepaspectratio]{fig3_dimensional_hierarchy.png}
\caption{Signal preservation without enrichment across dimensional expansion. The interval correlation $r = 0.8674$ (blue, left axis) is perfectly preserved from Golay(24D) through Barnes-Wall(256D, 512D, 1024D), while chord correlation (orange, right axis) remains flat at approximately $-0.09$. Higher dimensions add no discriminatory power.}\label{fig:dimensional_hierarchy}
\end{figure}

The Barnes-Wall lattices, by contrast, \emph{perfectly preserve} the Golay interval correlation at $r = 0.8674$ across all tested dimensions (256, 512, 1024). However, they provide \emph{zero enrichment} for chord differentiation: chord correlation remains at $r \approx -0.0894$ identically across all dimensions. The signal is preserved but never enriched. This ``frozen correlation'' phenomenon is explained by the fact that Barnes-Wall expansion is a linear projection from the Golay code's structure: it inherits the exact same distance distribution, merely scaling the distance values by a dimension-dependent factor.

\subsection{Phase VII: Full Seed Space and Leech Centroid Analysis}

Phase~VII confirmed that expanding from 4-bit to 12-bit Gray codes (using all 12 bits of the seed space) produces \emph{identical} codewords to the 4-bit version, because Gray(0--11) in 12 bits only activates bits 8--11. The Leech lattice centroid analysis showed that all 12 pitch clouds have centroids at the origin (zero vector), since the 128-point expansions are symmetric about the origin. This eliminates centroid-based metrics as a viable chord differentiation strategy.

% ═══════════════════════════════════════════════════════════════════════════
\section{Prime Number Structures}\label{sec:primes}
% ═══════════════════════════════════════════════════════════════════════════

\subsection{Phase VIII: The Mersenne/Fermat Duality}\label{sec:mf_duality}

Phase~VIII introduced a fundamentally different approach. Instead of encoding pitches through the Golay coding pipeline, we investigated whether prime number structures---Mersenne primes and Fermat primes in particular---naturally connect to musical intervals through modular arithmetic.

\begin{definition}[Mersenne/Fermat Residues mod 144]\label{def:mf}
For all Mersenne exponents $p \geq 5$: $(2^p - 1) \bmod 144 \in \{31, 127\}$. For all Fermat indices $k \geq 2$: $(2^{2^k} + 1) \bmod 144 \in \{17, 113\}$. The XOR identity $31 \oplus 127 = 17 \oplus 113 = 96 = 2^5 + 2^6$ holds for both families.
\end{definition}

This duality is illustrated in Figure~\ref{fig:mf_duality}. The four residues occupy specific positions on the mod-144 circle, with the Mersenne and Fermat families separated by the XOR bridge value $96 = 2/3 \times 144$. The binary representations reveal that the two families share bits $\{0, 4\}$ in common and differ only in bits $\{5, 6\}$ (the XOR bridge).

\begin{figure}[htbp]
\centering
\includegraphics[width=\columnwidth, height=0.35\textheight, keepaspectratio]{fig4_mersenne_fermat_duality.png}
\caption{The Mersenne/Fermat duality in mod-144 space. Fermat residues (17, 113, blue) and Mersenne residues (31, 127, orange) occupy complementary positions on the mod-144 circle. Dashed lines connect each family's pair, both yielding XOR $= 96$. The residues encode the fourth ($17 \equiv 5 \pmod{12}$) and fifth ($31 \equiv 7 \pmod{12}$).}\label{fig:mf_duality}
\end{figure}

The critical musical observation is that these residues encode the two most important intervals in Western harmony:
\begin{align}
17 &\equiv 5 \pmod{12} \quad \text{(the Perfect Fourth)} \label{eq:fermat_mod} \\
31 &\equiv 7 \pmod{12} \quad \text{(the Perfect Fifth)} \label{eq:mersenne_mod}
\end{align}
Furthermore, $113 \equiv 5 \pmod{12}$ and $127 \equiv 7 \pmod{12}$, so the pattern holds for all four residues. This is not a coincidence: the fifth ($7$ semitones) has the form $2^3 - 1$ (a Mersenne number), and the fourth ($5$ semitones) equals the Fermat prime $F_0 = 5$. Since $7 + 5 = 12$ (the octave), we obtain the decomposition shown in Figure~\ref{fig:12tet}:
\begin{equation}\label{eq:12tet}
\text{12-TET} = \underbrace{7}_{\text{Fifth} = 2^3 - 1\text{ (Mersenne form)}} + \underbrace{5}_{\text{Fourth} = F_0\text{ (Fermat prime)}}
\end{equation}

\begin{figure}[htbp]
\centering
\includegraphics[width=\columnwidth, height=0.35\textheight, keepaspectratio]{fig8_12tet_decomposition.png}
\caption{The 12-tone equal temperament octave decomposes into a Mersenne component (the fifth, 7 semitones $= 2^3 - 1$) and a Fermat component (the fourth, 5 semitones $= F_0$). Together they partition the chromatic circle.}\label{fig:12tet}
\end{figure}

\subsection{Pure Just-Intonation Exponent Vectors}

The most striking result of Phase~VIII came from abandoning the coding pipeline entirely. When pitches are encoded using the prime exponent vectors of their just-intonation frequency ratios---for example, the perfect fifth $3/2$ becomes $(e_2, e_3, e_5) = (-1, 1, 0)$, the major third $5/4$ becomes $(-2, 0, 1)$, and so on---the Pearson correlation with consonance rank reaches $r = +0.9613$ ($R^2 = 92.4\%$). This is the strongest interval correlation achieved in the entire study.

This result has a clear interpretation: consonance is fundamentally a property of the prime factorization of frequency ratios. Intervals with small prime exponents (perfect fifth, perfect fourth, major third) are consonant because their frequency ratios involve small primes raised to small powers. The Golay code correlation of $r = 0.8674$ captures an echo of this structure---the CoF ordering approximates the harmonic series---but the coding-theoretic layer dilutes the signal (passing JI exponents through Golay reduces $r$ from $0.96$ to $0.67$).

\subsection{Prime-Power Spectral Encoding}

An alternative prime-based approach computed spectral fingerprints by measuring each pitch's NRCI (Normalized Resonance Coherence Index) at power-of-prime frequencies ($p^{pc}$ for primes $p \in \{2, 3, 5, 7, 11, 13, 17, 19\}$). This produced the strongest chord signal observed up to that point: $r = -0.6185$ (inverted). The inversion indicates that \emph{dissonant} chords cluster together in spectral space while consonant chords spread apart---a structural separation principle, but in the wrong direction for a positive consonance metric.

% ═══════════════════════════════════════════════════════════════════════════
\section{Reverse-Engineering and the 4D Residue Fingerprint}\label{sec:phase9}
% ═══════════════════════════════════════════════════════════════════════════

\subsection{Phase IX: What Does Harmony Need?}

Phase~IX reversed the investigation's direction. Instead of asking ``can the UBP map harmonies?'' we asked ``what distance resolution does harmony need, and can any UBP component provide it?'' The requirements analysis established that harmonic differentiation requires at least six distinct distance values at the interval level (one per consonance rank) and at least four to five at the chord level (to separate major, minor, diminished, augmented, and dissonant triad types).

A Monte Carlo simulation confirmed that random 12-bit vectors provide an average of 8.7 distinct distances---more than sufficient. The bottleneck is the Golay code, which collapses these to three. Raw 12-bit Gray codes (without Golay encoding) produce four distinct distances $\{1, 2, 3, 4\}$ and achieve $r = +0.8017$ for intervals, confirming that the Golay code's distance amplification (from 4 to 2 effective levels) is responsible for both the strength and the ceiling of the $r = 0.8674$ result.

\subsection{The 4D Prime Residue Fingerprint}

The breakthrough came from representing each pitch class as a 4-dimensional vector of modular distances to the four prime residues:
\begin{equation}\label{eq:residue_vec}
\mathbf{v}_{pc} = \bigl(d(17, 2^{pc} \bmod 144),\; d(31, 2^{pc} \bmod 144),\; d(113, 2^{pc} \bmod 144),\; d(127, 2^{pc} \bmod 144)\bigr)
\end{equation}
where $d(a, b) = \min((a-b) \bmod 144, (b-a) \bmod 144)$ is the circular distance on $\mathbb{Z}/144\mathbb{Z}$.

For chords, the 4D vectors of the constituent pitches are aggregated (via mean distance, standard deviation, and spread). The \emph{standard deviation of the 4D residue distances} across a chord's pitches achieves:
\begin{itemize}[leftmargin=2em]
\item Interval correlation: $r = -0.50$ (Euclidean), moderate inverse
\item \textbf{Chord correlation: $r = -0.8790$ (Euclidean)} --- the best chord result in the entire 9-phase study
\end{itemize}

The mechanism is clear from the data: consonant chords (major triads, minor triads, sus chords) have pitches that \emph{spread} across the prime residue landscape, producing large 4D distance standard deviations (9.3--12.0). Dissonant clusters have pitches that are \emph{close} in residue space, producing small deviations (2.7--4.5). The correlation is inverted because consonant $\to$ spread (large values) but our ranking has consonant $\to$ small rank numbers.

% ═══════════════════════════════════════════════════════════════════════════
\section{The Prime-Layer Harmonic Module}\label{sec:phase10}
% ═══════════════════════════════════════════════════════════════════════════

\subsection{Phase X: Jaccard Similarity on Prime Residue Sets}

Phase~X addressed the inversion problem by switching from Euclidean distance to Jaccard similarity. The motivation is structural: if consonant chords spread across the prime residue landscape, then their pitch proximity sets should have \emph{less} overlap (lower Jaccard similarity) than dissonant clusters. But we want a \emph{positive} correlation (higher similarity $\to$ more consonant), so we construct proximity sets differently.

\subsubsection{Set Construction}

For each pitch class $pc$, we compute $2^{pc} \bmod 144$ and construct a binary proximity set:
\begin{equation}\label{eq:prox_set}
S_{pc}(\tau) = \bigl\{r \in \{17, 31, 113, 127\} : d(r, 2^{pc} \bmod 144) \leq \tau\bigr\}
\end{equation}
where $\tau$ is a proximity threshold. For a chord, the chord set is the union of its constituent pitch proximity sets.

\subsubsection{The Threshold Sweep}

The choice of threshold $\tau$ is critical. Figure~\ref{fig:jaccard_threshold} shows the chord correlation as a function of $\tau$ across the range $[1, 72]$. The correlation rises from zero at $\tau = 5$, peaks sharply at $\tau = 10$ with $r = +0.8244$, then declines. This peak is the first \emph{positive} chord differentiation above $r = 0.5$ achieved in the entire ten-phase study.

\begin{figure}[htbp]
\centering
\includegraphics[width=\columnwidth, height=0.35\textheight, keepaspectratio]{fig5_jaccard_threshold.png}
\caption{Chord correlation (Pearson $r$) versus Jaccard proximity threshold. The sharp peak at $\tau = 10$ yields $r = +0.8244$, the first positive chord differentiation above $r = 0.5$ in the study. The shaded region highlights the optimal zone.}\label{fig:jaccard_threshold}
\end{figure}

\subsubsection{Why Threshold 10? The Structural Explanation}

The threshold $\tau = 10$ is not arbitrary; it is determined by the geometry of the prime residues in mod-144 space. Residue $17$ (the Fermat anchor) sits at position 17, and pitch class 7 ($G$, the dominant) satisfies $d(17, 2^7 \bmod 144) = d(17, 128) = d(17, 128 - 144) = d(17, -16) = 33$, which modulo 144 circular distance is $\min(128 - 17, 144 - 128 + 17) = \min(111, 33) = 33$. The key boundary falls between pitch classes $F\#$ ($d = 11$) and $G$ ($d = 10$) from residue 17.

More fundamentally, the threshold $\tau = 10$ creates two pitch groups: $\{C, C\#, D, D\#, E, F\}$ (within distance 10 of residue 17) and $\{G, G\#, A, A\#, B\}$ (beyond). This corresponds to the tritone divide. Consonant triads systematically span this boundary, giving them lower average Jaccard similarity across pitch pairs. Dissonant clusters tend to stay within one group, yielding higher similarity. The sign flips because we now measure \emph{similarity}: consonant chords are \emph{less} similar in their prime residue proximity, producing a positive correlation when we invert the metric.

\subsection{Transposition Invariance}

A crucial validation is that the Jaccard encoding is perfectly transposition-invariant. All 12 major triads produce identical pairwise average Jaccard similarity ($J = 0.3333$), as do all 12 minor triads ($J = 0.3333$). This is not an artifact of small sample size but a structural consequence of the modulo arithmetic: transposing a chord by $k$ semitones adds $k$ to each pitch class, which shifts all $2^{pc} \bmod 144$ values by a multiplicative factor of $2^k$, preserving the relative distances to the four prime residues.

Furthermore, 12/12 major triads and 12/12 minor triads span the $\tau = 10$ boundary (i.e., their pitches fall in both proximity groups), while only 4/12 tested dissonant clusters do so. This is the mechanism: consonant triads \emph{must} span the boundary because they contain intervals (third and fifth) that cross it, while dissonant clusters \emph{can} avoid it.

\subsection{The Complete Prime-Layer Module}

The complete Prime-Layer Harmonic Module, summarized in Figure~\ref{fig:pipeline}, operates as follows:

\begin{enumerate}[leftmargin=2em]
\item \textbf{Pitch encoding:} Each pitch class $pc$ maps to a 4D vector of circular distances to residues $\{17, 31, 113, 127\}$ in $\mathbb{Z}/144\mathbb{Z}$.
\item \textbf{Interval metric:} Jaccard similarity on binary feature sets derived from the 4D vector achieves $r = +0.7563$ with 18 unique similarity values.
\item \textbf{Chord metric:} Jaccard similarity on proximity sets with threshold $\tau = 10$ achieves $r = +0.8244$.
\item \textbf{Independence:} This layer operates entirely independently of the Golay/Leech coding layer. The coding layer handles error correction; the prime layer handles harmony.
\end{enumerate}

\begin{figure}[htbp]
\centering
\includegraphics[width=\columnwidth, height=0.35\textheight, keepaspectratio]{fig7_prime_pipeline.png}
\caption{The Prime-Layer Harmonic Module pipeline. Pitch classes are mapped through $2^{pc} \bmod 144$ to proximity sets relative to the four prime residues, and Jaccard similarity provides both interval and chord consonance metrics.}\label{fig:pipeline}
\end{figure}

\subsection{Independent Confirmation via Mersenne Number Factorization}

An independent line of evidence corroborates the prime-layer approach. Instead of using residue proximity sets, we computed the Jaccard similarity of the actual \emph{prime factor sets} of the Mersenne numbers $2^{pc} - 1$ for each pitch class. This achieved chord correlation $r = +0.5382$ using genuine Mersenne number prime factors, not modular residues. While weaker than the proximity set method, this confirms that the harmonic signal is present in the prime factorization structure itself, not merely in the modular arithmetic artifact.

% ═══════════════════════════════════════════════════════════════════════════
\section{Discussion}\label{sec:discussion}
% ═══════════════════════════════════════════════════════════════════════════

\subsection{The Complete Consonance Signal Landscape}

Figure~\ref{fig:landscape} presents the complete landscape of consonance correlations across all methods tested in this ten-phase investigation. Several patterns emerge from this panoramic view.

\begin{figure}[htbp]
\centering
\includegraphics[width=\columnwidth, height=0.4\textheight, keepaspectratio]{fig6_consonance_signal_landscape.png}
\caption{The consonance signal landscape across all methods. Blue bars indicate positive correlations (consonance increases with the metric); red bars indicate inverted correlations (consonance decreases). Methods are grouped by whether they target intervals or chords.}\label{fig:landscape}
\end{figure}

First, the strongest interval signal ($r = +0.9613$) comes from pure JI exponent vectors, confirming that consonance is fundamentally a property of prime factorization. The coding-theoretic approaches (CoF Gray, Leech active dimension overlap) capture echoes of this signal but at lower correlation values and often with inversion.

Second, chord differentiation proved substantially harder than interval differentiation across every method tested. The best chord result ($r = +0.8244$, Jaccard proximity sets) required the prime-layer module; no purely coding-theoretic method exceeded $|r| = 0.3$ for chords.

Third, the distinction between positive and inverted correlations is methodologically important. Inverted correlations ($r < 0$) indicate that the metric captures \emph{some} harmonic structure but in the wrong direction---typically because dissonant entities cluster together (they share fewer prime factors or span less of the residue landscape). The Jaccard proximity method succeeds because it converts this spread-based separation into a positive similarity-based metric.

\subsection{Why $144 = 12^2$ Is the Correct Modulus}

The modulus $144 = 2^4 \times 3^2 = 12^2$ is not arbitrary. It is the smallest modulus that simultaneously:
\begin{enumerate}[leftmargin=2em]
\item Squares the octave (12), placing residues at octave-plus-interval positions (e.g., $17 = 12 + 5$ = ``one octave above, plus the fourth'');
\item Provides four distinct Mersenne/Fermat residues rather than two (mod~96 collapses to two residues);
\item Creates the $\tau = 10$ boundary at the tritone divide between $F\#$ and $G$, which is the most natural partition point in the chromatic system.
\end{enumerate}
Residue $17$ occupies position $17 = 12 + 5$, which is literally ``one octave above C, plus the fourth.'' Residue $31$ sits at $31 = 2 \times 12 + 7$, ``two octaves above C, plus the fifth.'' This octave-stacking is what makes the residues musically meaningful: they encode intervals \emph{relative to octaves}, not just as bare numbers.

\subsection{The Sporadic Groups Connection}

A curious observation from Phase~VIII is that all 26 sporadic simple groups have orders divisible by 12 (i.e., $|G| \equiv 0 \pmod{12}$). While we do not claim a causal connection to music theory, this divisibility property means that the sporadic groups ``live in the octave'' in a precise number-theoretic sense. The Mathieu group $M_{12}$, which automorphizes the Steiner system $S(5, 6, 12)$ (132 hexachords on 12 points), has a particularly direct connection to the 12-pitch-class system.

\subsection{Limitations and Future Work}

Several limitations should be acknowledged. First, the consonance ranking used here is a simplified six-level scale; a continuous perceptual consonance measure would provide finer resolution. Second, the Jaccard threshold $\tau = 10$ was identified post hoc through a sweep; a theoretical derivation of the optimal threshold from first principles would strengthen the result. Third, the study considers only equal-tempered pitch classes; extending to just intonation or microtonal systems would test the generality of the prime-layer approach. Fourth, the chord set (20 chords) is relatively small; larger chord sets including extended harmonies (9ths, 11ths, 13ths) would provide a more rigorous test.

% ═══════════════════════════════════════════════════════════════════════════
\section{Conclusion}\label{sec:conclusion}
% ═══════════════════════════════════════════════════════════════════════════

This ten-phase computational investigation has yielded both definitive negative results and constructive positive results. On the negative side, we have proven that the extended binary Golay code $[24, 12, 8]$, the Leech lattice $\Lambda_{24}$, and the Barnes-Wall lattice family are fundamentally incapable of chord-level harmonic differentiation. The Golay code's three-distance constraint ($\{8, 12, 16\}$) is an absolute ceiling: no UBP metric defined on these structures can exceed $|r| \approx 0.3$ for chords. Dimensional expansion preserves but never enriches the interval signal.

On the positive side, we have demonstrated that the consonance signal resides in prime number theory, not in coding theory. Pure JI exponent vectors achieve $r = +0.9613$ for intervals, and a Jaccard-based Prime-Layer Harmonic Module operating on Mersenne/Fermat residue proximity sets achieves $r = +0.8244$ for chords. The mechanism is structurally explained by the fact that the four prime residues $\{17, 31, 113, 127\}$ in mod-144 space encode the fourth and fifth---the two intervals that partition the octave---and that consonant triads systematically span the resulting boundary while dissonant clusters do not.

The decomposition $12\text{-TET} = (2^3 - 1) + F_0 = 7 + 5$ (Mersenne plus Fermat) provides a satisfying number-theoretic explanation for why 12-tone equal temperament works as well as it does: the octave naturally decomposes into a Mersenne number and a Fermat prime, the two most important families in number theory. This is not a design choice but a mathematical fact, and the Prime-Layer Harmonic Module presented here operationalizes this fact into a working consonance metric.

The broader methodological lesson is that ``honest mapping''---the practice of following the mathematics wherever it leads, including through negative results---ultimately produces stronger results than forcing a predetermined framework onto the data. The coding-theoretic approach failed for chords; admitting this failure led us to number theory, which succeeded. The UBP system, designed for error correction, does not naturally model harmony. But the prime number structures that the UBP's prime-testing subsystem brushes against---Mersenne and Fermat primes in particular---do.

% ═══════════════════════════════════════════════════════════════════════════
% REFERENCES
% ═══════════════════════════════════════════════════════════════════════════

\begin{thebibliography}{99}

\bibitem{golay1949}
M.~J.~E.~Golay,
``Notes on digital coding,''
\textit{Proceedings of the IRE}, vol.~37, no.~6, p.~657, 1949.

\bibitem{conway1985}
J.~H.~Conway, R.~T.~Curtis, S.~P.~Norton, R.~A.~Parker, and R.~A.~Wilson,
\textit{Atlas of Finite Groups: Maximal Subgroups and Ordinary Characters for Simple Groups}.
Oxford University Press, 1985.

\bibitem{leech1967}
J.~Leech,
``Notes on sphere packings,''
\textit{Canadian Journal of Mathematics}, vol.~19, pp.~251--267, 1967.

\bibitem{sloane1977}
N.~J.~A.~Sloane,
``Error-correcting codes and invariant theory: Applications to binary Golay codes,''
\textit{SIAM Journal on Algebraic Discrete Methods}, vol.~3, no.~2, pp.~176--192, 1982.

\bibitem{plomp1965}
R.~Plomp and W.~J.~M.~Levelt,
``Tonal consonance and critical bandwidth,''
\textit{Journal of the Acoustical Society of America}, vol.~38, no.~4, pp.~548--560, 1965.

\bibitem{sethares1993}
W.~A.~Sethares,
``Local consonance and the relationship between timbre and scale,''
\textit{Journal of the Acoustical Society of America}, vol.~94, no.~4, pp.~2118--2128, 1993.

\bibitem{tymoczko2006}
D.~Tymoczko,
``The geometry of musical chords,''
\textit{Science}, vol.~313, no.~5783, pp.~72--74, 2006.

\bibitem{ammo2011}
A.~Ammoch and T.~Noll,
``Fermat primes, Mersenne primes, and the musical scale,''
in \textit{Mathematics and Computation in Music}, Springer, 2011, pp.~1--12.

\bibitem{barnes1959}
E.~S.~Barnes and G.~M.~Wall,
``Some extreme forms defined in terms of abelian groups,''
\textit{Proceedings of the London Mathematical Society}, vol.~s3-9, no.~4, pp.~589--610, 1959.

\bibitem{griess1982}
R.~L.~Griess,
``The friendly giant,''
\textit{Inventiones Mathematicae}, vol.~69, no.~1, pp.~1--102, 1982.

\end{thebibliography}

\end{document}